\chapter{Background}

In this chapter, we cover some Definitions relevant for this thesis.

\section{Deterministic Finite Automaton}
\textbf{Definition:} A deterministic finite automaton (DFA) is a 5-tuple $M = \langle Q, \Sigma, \delta, q_0, F \rangle$ where
\begin{itemize}
    \item $Q$ is the finite set of states
    \item $\Sigma$ is the input alphabet
    \item $\delta : Q \times \Sigma \rightarrow Q$ is the transition function
    \item $q_0 \in Q$ is the start state
    \item $F \subseteq Q$ is the set of accept states (or final states).
\end{itemize}
For readability we define the product $q \cdot a = \delta(q, a)$, where $q \in Q$ is a state and $a \in \Sigma$ is a symbol of the alphabet.
Additionally we define $q \cdot w = (q \cdot a) \cdot x$ where $w = ax \in \Sigma^*$, $a \in \Sigma$ and $x \in \Sigma^*$.

\textbf{Definition:} DFA $M = \langle Q, \Sigma, \delta, q_0, F \rangle$ accepts the word $w \in \Sigma^*$ if $q \cdot w \in F$.

\textbf{Definition:} Let $M$ be a DFA. The language recognized by $M$ is defined as $\mathcal L(M) = \{w \in \Sigma^* \mid w \text{ is accepted by } M\}$.

We depict a DFA visually as a directed graph with labelled arcs, where there is a corresponding node for each state and a labelled arc for each transition in the transition function.

